\chapter{Implementation}
\label{chap:implementation}

\section{Repository Structure}

The implementation is contained in the benchmarking framework repository. The core modules are:
\begin{itemize}
    \item \code{benchmarking/core/config.py}: typed workload and benchmark configuration objects, including configuration hashing;
    \item \code{benchmarking/core/engine.py}: the common engine interface;
    \item \code{benchmarking/core/result.py}: benchmark result data structures;
    \item \code{benchmarking/runner/runner.py}: warm-up, repeated timing and aggregation;
    \item \code{benchmarking/storage/database.py}: SQLite persistence and schema handling;
    \item \code{benchmarking/workloads/mc_cpu.py}: NumPy CPU workloads;
    \item \code{benchmarking/workloads/mc_jax.py}: JAX workloads and AD support;
    \item \code{benchmarking/workloads/mc_cpp.py}: Python bindings for the C++ engine;
    \item \code{benchmarking/workloads/mc_rust.py}: Python bindings for the Rust engine;
    \item \code{benchmarking/cloud/metadata.py} and \code{benchmarking/cloud/pricing.py}: cloud metadata and cost support; and
    \item \code{experiments/run_profiler_vs_grid.py}: the final profiler-versus-grid experiment driver.
\end{itemize}

The C++ implementation is under \code{benchmarking/cpp/}, and the Rust implementation is under \code{benchmarking/rust/}. The report artefact generator is \code{Final_Year_Project_Report/evaluation/generate_report_assets.py}; it reads the stored results and writes the tables and figures used in Chapter~\ref{chap:evaluation}.

\section{Configuration and Workload Registration}

Workloads are configured with explicit fields for option parameters, path count, number of time steps, engine, AD mode and random seed. Configurations are hashed so that repeated runs can be identified and compared. This is useful when the same logical benchmark is run on different machines or with different profiler decisions.

The implementation treats workload support as a first-class concern. An engine can reject unsupported combinations, for example when AD is requested for an engine that does not implement it. This avoids accidental comparisons between incomplete or silently degraded implementations.

\section{Engine Implementations}

\subsection{NumPy CPU}

The NumPy engine provides a readable vectorised baseline. It is useful both as a reference implementation and as a lower-complexity comparison point for compiled engines. It supports the evaluated workloads without AD.

\subsection{JAX/XLA}

The JAX engine implements the workloads in a form suitable for JIT compilation and AD. Forward- and reverse-mode sensitivities are measured for supported workloads. JAX compilation effects are handled by warm-up before repeated timing, so the reported runtime is intended to measure steady-state execution rather than first-call compilation.

\subsection{C++/OpenMP}

The C++ engine is exposed to Python through a native extension. It is used for no-AD timing comparisons and provides a compiled implementation with explicit parallelism. The C++ engine does not provide the AD measurements reported in this project.

\subsection{Rust/Rayon}

The Rust engine is exposed through Python bindings and uses Rayon-style data parallelism for no-AD workloads. It provides the fastest measured no-AD European runtime in the final cloud experiment and is also strong on the local-volatility workload. As with C++, it is evaluated as a no-AD engine.

\subsection{CUDA Scaffold}

The repository contains CUDA-related workload scaffolding, but GPU execution was not included in the final evaluated experiments. The report therefore does not claim GPU results.

\section{Benchmark Runner}

The runner is responsible for applying a consistent measurement protocol. Each benchmark has warm-up runs followed by repeated timed runs. The runner records runtime summaries and derived throughput, and it forwards workload outputs such as prices and Greeks into the common result structure. Separating this logic from the engines keeps measurement policy consistent across implementations.

For AD overhead, the runner compares JAX AD timings against the corresponding JAX no-AD baseline for the same workload and path count. This keeps the overhead measurement focused on the cost of derivative computation within the same runtime system.

\section{SQLite Storage}

Results are written to \code{results/benchmarks.db}. SQLite was sufficient for this project because experiments are moderate in size, local inspection is convenient, and SQL queries can directly produce analysis slices. The schema includes fields for runtime, throughput, price, Greeks, analytical errors, experiment identifiers, cloud metadata, profiler decisions and SHA-specific metadata.

The storage layer enables a useful workflow: experiment scripts write raw results once, and report tables or figures can be regenerated later without rerunning benchmarks. This reduces the risk of transcription errors and supports the requirement that report claims be backed by actual result files.

\section{Cloud and Cost Metadata}

Cloud experiments attach instance type, region, zone and pricing information to each run. The final headline experiment used \code{n2-standard-4} in \code{europe-west1-b} at git commit \code{13789b8}. Cost per run is computed from runtime and hourly instance price. Because cloud prices and machine availability can change, the stored values are treated as experiment metadata rather than timeless constants.

\section{Profiler Implementation}

The profiler experiment driver constructs the full candidate grid, executes probe runs, applies old-profiler and SHA selection logic, and records which configurations would be fully benchmarked by each method. For the final experiment, the full-grid baseline consisted of 16 configurations evaluated at three full path counts, for 48 full benchmark cells. The old profiler selected 30 full cells, while SHA selected 24.

The SHA implementation records intermediate promotion decisions, including active configurations at each probe round. These records are used to generate the SHA progression plot in Chapter~\ref{chap:evaluation}. The implementation also stores enough information to compare selected sets with the full grid using Pareto recovery and regret metrics.

\section{Testing}

The repository includes system and AD-focused tests in \code{tests/test_system.py} and \code{tests/test_ad_framework.py}. These tests cover configuration construction, engine interfaces, result storage and AD-related behaviour. They do not replace full numerical validation, but they help protect the framework mechanics that the experiments rely on.
